\paragraph{单位与\trQuasiequivalence。}\label{sub:units-quasi-equivalences}
若 $\scr{A}$ 是\trStrictlyUnital{}的 \trACategory，那么 $Tw\scr{A}$ 亦然。具体说来，%
\trTwistedComplex{} \eqref{eq:pretwisted-complex} 的严格\trIdentityMorphism{}即“单位矩阵”%
（identity matrix）$(\phi^{ji}\otimes x^{ji})$，其中
$\phi^{ii}\otimes x^{ii}=\id_{V^i}\otimes e_{X^i}$，而所有非对角元均为零。此外，若
$\scr{G}\colon\scr{A}\to\scr{B}$ 是\trStrictlyUnital{}的 \trAFunctor，那么 $Tw\scr{G}$
亦然。最后，给定\trStrictlyUnital{}的 $\scr{A}$ 与任意的 $\scr{C}$，
\begin{equation}\label{eq:tw-nufun-functor}
Tw\colon\mathit{nu}\text{-}\mathit{fun}(\scr{C},\scr{A})
\longrightarrow\mathit{nu}\text{-}\mathit{fun}(Tw\scr{C},Tw\scr{A})
\end{equation}
是一个\trStrictlyUnital{}的 \trAFunctor。要看出这一点，只需回忆\cref{sub:functor-categories}
中 $E_{\scr{G}}$ 的表达式，并把它代入\eqref{eq:tw1-map}。较不显然的\trCunital{}类比则是：

\vkthmsync{24}
\begin{lemma}\label{thm:tw-preserves-cunital}
(i) 若 $\scr{A}$ 是\trCunital{}的，那么 $Tw\scr{A}$ 亦然。
(ii) 设 $\scr{A}$ 与 $\scr{B}$ 都是\trCunital{}的，且 $\scr{G}\colon\scr{A}\to\scr{B}$
是一个\trCunital{}的 \trAFunctor。那么 $Tw\scr{G}$ 亦然。
(iii) 对\trCunital{}的 $\scr{A}$ 与任意的 $\scr{C}$，\eqref{eq:tw-nufun-functor}
是\trCunital{}的 \trAFunctor。
\end{lemma}
\begin{proof}
(i) 我们使用严格化论证，一如\cref{thm:functor-category-cunital}的证明中那样。具体说来，
取一个\trFormalDiffeomorphism{} $\Phi$，使得 $\tilde{\scr{A}}=\Phi_{*}\scr{A}$ 是%
\trStrictlyUnital{}的。由\cref{thm:tw-preserves-cohomologically-ff}，诱导出的%
\trNonunital{} \trAFunctor{} $Tw\Phi\colon Tw\scr{A}\to Tw\tilde{\scr{A}}$
是\trCohomology\trFullyFaithful{}的。既然后一个\trCategory{}是\trStrictlyUnital{}的，
前一个就必定是\trCunital{}的。
(ii) 这次考虑 $\tilde{\scr{G}}=\scr{G}\circ\Phi^{-1}\colon\tilde{\scr{A}}\to\scr{B}$，
其中 $\Phi^{-1}$ 是逆\trFormalDiffeomorphism。取某个 $X\in Tw\tilde{\scr{A}}$，并为其
配备适当的\trFiltration{} $F^kX$，再设 $e_X$ 是它在 $Tw\tilde{\scr{A}}$ 中的严格%
\trIdentityMorphism。按定义，像 $(Tw\tilde{\scr{G}})^1(e_X)$ 相对于 $Tw\tilde{\scr{G}}(X)$
的像\trFiltration{} $Tw\tilde{\scr{G}}(F^kX)$ 是下对角（lower-diagonal）的，而在诸商
$Tw\tilde{\scr{G}}(F^kX)/Tw\tilde{\scr{G}}(F^{k+1}X)$ 上诱导的\trEndomorphism{}都是代表%
\trIdentityMorphism{}的链。利用\eqref{eq:twisted-spectral-sequence} 便可看出，这蕴含：
与 $[Tw\tilde{\scr{G}}^1(e_X)]$ 的左乘和右乘都是
$\Hom_{H(Tw\scr{B})}(Tw\tilde{\scr{G}}(X),Tw\tilde{\scr{G}}(X))$ 的\trAutomorphism，从而
$[Tw\tilde{\scr{G}}^1(e_X)]$ 本身是一个\trInvertible{}元。由 (i) 我们已知 $Tw\Phi$
是\trCunital{}的，所欲证的命题便由此而来。

(iii) 取定一个\trNonunital{}的 \trAFunctor{} $\scr{G}\colon\scr{C}\to\scr{A}$，并令
$\tilde{\scr{G}}=\Phi\circ\scr{G}$，其中 $\Phi$ 如前。作为 \eqref{eq:tw-left-right-composition}
的特例，我们有交换图
{\footnotesize\begin{tikzcd}[column sep=tiny,nodes={inner sep=1pt}]
\hom_{\mathit{nu}\text{-}\mathit{fun}(\scr{C},\scr{A})}(\scr{G},\scr{G})
  \ar[r,"\scr{L}^1_{\Phi}"]\ar[d,"Tw^1"']
& \hom_{\mathit{nu}\text{-}\mathit{fun}(\scr{C},\tilde{\scr{A}})}
  (\tilde{\scr{G}},\tilde{\scr{G}})\ar[d,"Tw^1"]\\
\hom_{\mathit{nu}\text{-}\mathit{fun}(Tw\scr{C},Tw\scr{A})}(Tw\scr{G},Tw\scr{G})
& \hom_{\mathit{nu}\text{-}\mathit{fun}(Tw\scr{C},Tw\tilde{\scr{A}})}
  (Tw\tilde{\scr{G}},Tw\tilde{\scr{G}}).\ar[l,"\scr{L}^1_{Tw\Phi^{-1}}"{swap,pos=0.3}]
\end{tikzcd}}
由上面的 (ii) 与\cref{thm:left-composition-cunital}可见，在\trCohomology{}层面上，
两条\trHorizontal{}箭头把恒等映到恒等。右边的\trVertical{}箭头亦是如此（那里我们有
严格单位），于是交换性便蕴含了所欲证的命题。
\end{proof}

\begin{lemma}\label{thm:tw-quasi-equivalence-preserved}
设 $\scr{G}\colon\scr{A}\to\scr{B}$ 是\trCunital{}的 \trACategory{}之间的一个%
\trQuasiequivalence。那么 $Tw\scr{G}\colon Tw\scr{A}\to Tw\scr{B}$ 也是一个%
\trQuasiequivalence。
\end{lemma}
这是\cref{thm:quasi-equivalence-homotopy-inverse}与前一条引理的直接推论：先在相应的%
\trFunctor\trCategory{}中把 $\scr{G}$ 在\trIsomorphism{}的意义下求逆，再考虑诱导的%
\trFunctor{}与\trNaturalTransformation{}即可（另一种做法不那么形式化，而更多地依赖%
\trExactTriangle，见下文\cref{thm:tw-yoneda-combination-cff}的证明）。
